\documentclass{book}

\usepackage{geometry}
\usepackage{amsmath}
\usepackage{amssymb}
\usepackage{amsthm}
\usepackage{amsmath}
\usepackage{tgpagella}

\theoremstyle{definition}
\newtheorem{definition}{Definition}[section]

\begin{document}

\chapter{Simulations}

\section{Problem definitions}

\subsection{Non-diffusing receptors on the cell boundary}

This problem assumes the receptor lives on the cell boundaries without diffusing. The ligand is able to diffuse across the entire extracellular space, and interacts with the receptor at the cell boundaries according to the relevant ordinary differential equation. Production of both ligand and receptor exclusively takes place on the cell boundaries.

Let $\Omega \subseteq \mathbb{R}^{m}$ be the open and bounded domain and denote the cells by an open set $C \subseteq \Omega$ of which the closure $\overline{C}$ is contained in $\Omega$. Furthermore, let $E = \Omega \setminus \overline{C}$ be the extracellular space, and let $l(t, x)$ denote the concentration of ligands at time $t > 0$ and point $x \in \overline{E}$, and $r(t, x)$ the concentration of receptors at time $t > 0$ and point $x \in \partial C$. Then the unknowns $l$ and $r$ should satisfy

\begin{equation}
    \begin{aligned}
        r_t &= \gamma (a - r + r^{2}l) && \text{on } (0, \infty) \times \partial C, \\
        l_t &= d(???) \Delta l && \text{in } (0, \infty) \times E, \\
        \frac{\partial l}{\partial \nu} &= \gamma (b - r^{2}l) && \text{on } (0, \infty) \times \partial C, \\
        \frac{\partial l}{\partial \nu} &= 0 && \text{on } (0, \infty) \times \partial \Omega, \\
    \end{aligned}
\end{equation}

\noindent
where $d > 0$ is the diffusion coefficient, $\gamma > 0$ is the reaction coefficient and $a > 0$ and $b > 0$ are the production coefficients of the receptor and ligand respectively.

Let $\tau > 0$ be the time discretization parameter and for every $n \in \mathbb{N}$, let $t_n := n\tau$ and

\begin{equation}
    r_n(x) \approx r(t_n, x) = r(n\tau, x) \quad \text{and} \quad l_n(x) \approx l(t_n, x) = l(n\tau, x)
\end{equation}

\noindent
be the approximations of receptor and ligand at time $t_n = n\tau$ respectively. Since we only consider an ODE for $r$, we do not need to formulate a variational form for this unknown. Instead, we will approximate the ODE for $r$ with some finite difference quotient (???),

\begin{equation}
        \frac{r_{n + 1} - r_{n}}{\tau} = \gamma (a - r_{n} + r_{n}^{2}l_{n}) \quad \text{on } \partial C, n \in \mathbb{N},
\end{equation}

\noindent
using forward Euler. This gives us an explicit formula for $r_{n + 1}$,

\begin{equation}
        r_{n + 1} = r_{n} + \tau \gamma (a - r_{n} + r_{n}^{2}l_{n}) \quad \text{on } \partial C, n \in \mathbb{N},
\end{equation}

\noindent
with $r_{0}$ initialized as the exact discretization of the initial conditions on $r$. To approximate $l$, we first discretize time using a method slightly different as the one used for $r$,

\begin{equation}
        \frac{l_{n + 1} - l_{n}}{\tau} = d \Delta l_{n + 1} \quad \text{in } E, n \in \mathbb{N}.
\end{equation}

\noindent
The use of the backward Euler instead of forward Euler method provides additional stability, which is particularly convenient considering we are dealing with a second-order derivative. This gives us the implicit formula for $l_{n + 1}$,

\begin{equation}
        l_{n + 1} - \tau d \Delta l_{n + 1} = l_{n} \quad \text{in } E, n \in \mathbb{N}.
\end{equation}

\noindent
As for $r$, $l_{0}$ is initialized as the exact discretization of the initial condition on $l$. To derive the variational form of this problem, we multiply both sides by a test function $\phi$, integrate on $E$ and use partial integration,

\begin{equation}
    \begin{aligned}
        \int_{E} l_n &= \int_{E} \left[l_{n + 1} \phi - \tau d \Delta l_{n + 1} \phi \right] \\
        &= \int_{E} l_{n + 1} \phi + \tau d \int_{E} \nabla l_{n + 1} \nabla \phi - \tau d \int_{\partial E} \frac{\partial l_{n + 1}}{\partial \nu} \phi\ dS \\
        &= \int_{E} l_{n + 1} \phi + \tau d \int_{E} \nabla l_{n + 1} \nabla \phi - \tau d \int_{\partial C} \gamma (b - r_{n + 1}^{2}l_{n + 1}) \phi\ dS, \\
    \end{aligned}
\end{equation}

\noindent
where in the last step we used $\partial E = \partial \Omega\ \cup\ \partial C$, since $\overline{C} \subseteq \Omega$, and the Neumann boundary conditions.

\end{document}